\documentclass[12pt]{article}





\usepackage{amsmath}  % For mathematical symbols and equations
\usepackage{amssymb}  % For additional math symbols
\usepackage{geometry} % To control page margins
\geometry{a4paper, margin=1in}





\usepackage{hyperref} % To include hyperlinks
\usepackage{fancyhdr} % To customize the header and footer
\pagestyle{fancy}
\fancyhf{}
\fancyhead[L]{Programming Paradigms} 
\fancyfoot[C]{\thepage}





\title{Names,Bindings and Scope}
\author{Saad Almalki}
\date{\today}

\begin{document}

\maketitle

\section{Introduction} 
In an imperative language, a variable is an abstraction of a memory cell

\section{Binding}
Binding is associate entity with attribute.

\subsection{Possible Binding Times}

\begin{itemize}
    \item \textbf{Language design time}
    \item \textbf{Implementation time}
    \item \textbf{Compile time} : Bindings a variable with a type
    \item \textbf{Load time} : Bindings a static local variables
    \item \textbf{Run time} : Bindings a nonstatic local variables
\end{itemize}


\section{Names} 

\subsection{Length}
If They too short , they cannot connotative.

\subsubsection{Language examples}
\begin{itemize}
    \item FORTRAN I : Maximum 6
    \item COBOL : Maximum 30
    \item FORTRAN 90 and C 89 : Maximum 31
    \item C 99 : Maximum 63
    \item Java,C++,Csharp : No limits
\end{itemize}


\subsection{Case sensitivity}

C-Based language is case sensitive but others are not.

\section{Special words}

\subsection{Keywords}
word is used in Certain contexts

\subsection{Reserved words}
Words cannot be used as pre-defined words , FORTRAN have more than 300 reserved words !!!

\subsection{Special characters}

- PHP: all variable names must begin with Dollar signs

– Perl: all variable names begin with special characters, which specify the variable’s type
-@ are Ruby: variable names that begin with `@@` are instance variables

examples : if , else , int

\section{Attributes}

Variables can be characterized as a sextuple of attributes:

\begin{itemize}
    \item \textbf{Name}
    \item \textbf{Value}
    \item \textbf{Address}
    \item \textbf{Lifetime}
    \item \textbf{Scope}
    \item \textbf{Type}
    \end{itemize}


\section{Declaration}
two kinds of declarations

\begin{itemize}
    \item \textbf{Explicit declaration} : int x = 5;
    \item \textbf{Implicit declaration} : x = 5;
\end{itemize}

\subsection{Type binding}
When? - three kinds of type bindings:

1. Static type binding 2. Dynamic type binding 3. Type inference



\begin{flushright}
Created by: Saad Almalki \\
\end{flushright}

\end{document}
